\documentclass[12pt]{article}
\usepackage[margin=1in]{geometry}
\usepackage{amsmath}
\usepackage{graphicx}
\usepackage{siunitx}
\usepackage{booktabs}
\usepackage{float}
\usepackage{hyperref}

\begin{document}

\begin{titlepage}
    \centering
    \vspace*{1in}
    {\Large\bfseries University of South Carolina \par}
    \vspace{1.5cm}
    {\Large EMCH 362: Lab 2 Heat Transfer \par}
    \vspace{2cm}
    {\huge\bfseries Determination of Thermal Conductivity of Brass \par}
    \vspace{2cm}
    {\Large Authors: \par}
    \vspace{0.5cm}
    {\large Jacob Vaught \par}
    {\large Michael Barger \par}
    {\large Athena Shier \par}
    \vfill
    {\large Principal Investigators: \par}
    \vspace{0.5cm}
    {\large Mario Villegas \par}
    {\large Mobli Mostafa \par}
    \vspace{1cm}
    {\large \today \par}
\end{titlepage}

\pagenumbering{arabic}

\begin{abstract}
This experiment aimed to determine the thermal conductivity of a brass sample using a steady-state heat transfer method. A brass specimen was placed between a heat source and a heat sink, with temperatures measured at various points using thermocouples. The thermal conductivity was calculated using Fourier's law of conduction, simplified under steady-state and one-dimensional assumptions. The average thermal conductivity obtained was \num{121.06}\,\si{\watt\per\meter\per\celsius} with a standard deviation of \num{12.41}\,\si{\watt\per\meter\per\celsius}. This value aligns closely with literature values for brass, validating the experimental approach. Sources of error were analyzed, and the results demonstrate the effectiveness of the method in measuring thermal properties of materials.
\end{abstract}

\section{Introduction}

Thermal conductivity is a fundamental property that quantifies a material's ability to conduct heat. It plays a crucial role in engineering applications such as heat exchangers, thermal insulation, and electronic devices. Brass, an alloy of copper and zinc, is widely used due to its favorable thermal and mechanical properties. Accurate determination of its thermal conductivity is essential for efficient thermal management in various applications.

This experiment utilizes Fourier's law of conduction to measure the thermal conductivity of a brass sample under steady-state conditions. By simplifying the law with assumptions of one-dimensional heat flow and constant thermal conductivity, the experiment aims to calculate the thermal conductivity and compare it with established literature values.

\section{Methods}

\subsection{Experimental Setup}

A brass sample was prepared with the following dimensions (see Table~\ref{tab:sample_properties}).

\begin{table}[H]
    \centering
    \caption{Brass Sample Properties}
    \label{tab:sample_properties}
    \begin{tabular}{lc}
        \toprule
        Property & Value \\
        \midrule
        Diameter, $D$ & \SI{25}{\milli\meter} \\
        Length, $L$ & \SI{30}{\milli\meter} \\
        Cross-sectional Area, $A$ & \SI{4.9087e-4}{\square\meter} \\
        \bottomrule
    \end{tabular}
\end{table}

The sample was sandwiched between two identical hot layers and a cooling layer. Conducting grease was applied between all surfaces to ensure good thermal contact. The lateral sides of the sample were insulated to promote one-dimensional heat transfer. Eight thermocouples (AI0 to AI7) were embedded at specified distances along the heat flow path to measure temperature gradients.

\subsection{Procedure}

\begin{enumerate}
    \item Connected the power supply to the heat source and ensured proper electrical connections.
    \item Placed the brass specimen between the heat source and heat sink.
    \item Positioned the eight thermocouples at distances of \SIlist{0;15;30;75;90;135;150;165}{\milli\meter} from the heat source.
    \item Measured the vertical distances between thermocouples and calculated the cross-sectional area $A$ of the brass sample (see Table~\ref{tab:sample_properties}).
    \item Turned on the cooling water to establish a temperature gradient.
    \item Monitored the thermocouple readings using DAQami software until steady-state conditions were achieved.
    \item Recorded the voltage ($V = \SI{5}{\volt}$) and current ($I = \SI{0.78}{\ampere}$) to calculate the heat transfer rate (see Table~\ref{tab:electrical_parameters}).
    \item Collected temperature data over time to verify steady-state conditions.
\end{enumerate}

\begin{table}[H]
    \centering
    \caption{Electrical Parameters}
    \label{tab:electrical_parameters}
    \begin{tabular}{lc}
        \toprule
        Parameter & Value \\
        \midrule
        Voltage, $V$ & \SI{5}{\volt} \\
        Current, $I$ & \SI{0.78}{\ampere} \\
        Heat Transfer Rate, $q = V \times I$ & \SI{3.9}{\watt} \\
        \bottomrule
    \end{tabular}
\end{table}

\subsection{Calculations}

Under the assumptions, Fourier's law simplifies to:
\begin{equation}
    q = -kA \frac{\Delta T}{\Delta y}
\end{equation}
Rearranged to solve for thermal conductivity $k$:
\begin{equation}
    k = \frac{q \Delta y}{A \Delta T}
\end{equation}
Calculated $k$ for each segment between thermocouples to obtain multiple conductivity values for statistical analysis.

\section{Results and Discussion}

\subsection{Steady-State Verification}

The temperature readings at each thermocouple stabilized over time, indicating steady-state conditions. The negligible fluctuations in temperature confirmed the validity of the steady-state assumption.

\subsection{Temperature Profile}

A linear temperature gradient was observed when plotting temperature against distance (Figure~\ref{fig:temp_profile}), consistent with one-dimensional heat conduction. Minor deviations from linearity can be attributed to experimental uncertainties and imperfect insulation.

\begin{figure}[H]
    \centering
    \includegraphics[width=0.8\textwidth]{temperature_profile.png}
    \caption{Temperature vs. Distance from the Hot Side}
    \label{fig:temp_profile}
\end{figure}

\subsection{Thermal Conductivity Calculations}

The average temperatures at each thermocouple were calculated (see Table~\ref{tab:avg_temp}).

\begin{table}[H]
    \centering
    \caption{Average Temperatures at Thermocouples}
    \label{tab:avg_temp}
    \begin{tabular}{ccc}
        \toprule
        Thermocouple & Distance (mm) & Temperature (\si{\celsius}) \\
        \midrule
        AI0 & 0   & 39.17 \\
        AI1 & 15  & 38.18 \\
        AI2 & 30  & 37.02 \\
        AI3 & 75  & 33.84 \\
        AI4 & 90  & 32.83 \\
        AI5 & 135 & 29.91 \\
        AI6 & 150 & 28.99 \\
        AI7 & 165 & 28.15 \\
        \bottomrule
    \end{tabular}
\end{table}

Using these values, thermal conductivity $k$ was calculated for each segment (see Table~\ref{tab:k_values}).

\begin{table}[H]
    \centering
    \caption{Thermal Conductivity Calculations}
    \label{tab:k_values}
    \begin{tabular}{ccccc}
        \toprule
        Segment & $\Delta y$ (m) & $\Delta T$ (\si{\celsius}) & $k$ (\si{\watt\per\meter\per\celsius}) & Heat Flux $\left( \frac{q}{A} \right)$ (\si{\watt\per\meter\squared}) \\
        \midrule
        AI0 - AI1 & 0.015 & 0.99  & 120.22 & 7,944.5 \\
        AI1 - AI2 & 0.015 & 1.16  & 102.75 & 7,944.5 \\
        AI2 - AI3 & 0.045 & 3.18  & 112.44 & 7,944.5 \\
        AI3 - AI4 & 0.015 & 1.01  & 118.01 & 7,944.5 \\
        AI4 - AI5 & 0.045 & 2.92  & 122.60 & 7,944.5 \\
        AI5 - AI6 & 0.015 & 0.92  & 129.52 & 7,944.5 \\
        AI6 - AI7 & 0.015 & 0.84  & 141.86 & 7,944.5 \\
        \bottomrule
    \end{tabular}
\end{table}

The average thermal conductivity and standard deviation are summarized in Table~\ref{tab:summary_statistics}.

\begin{table}[H]
    \centering
    \caption{Summary of Thermal Conductivity Results}
    \label{tab:summary_statistics}
    \begin{tabular}{lc}
        \toprule
        Parameter & Value (\si{\watt\per\meter\per\celsius}) \\
        \midrule
        Average Thermal Conductivity, $\bar{k}$ & 121.06 \\
        Standard Deviation, $\sigma$ & 12.41 \\
        \bottomrule
    \end{tabular}
\end{table}

\subsection{Comparison with Literature}

Literature values for the thermal conductivity of brass range from \SIrange{100}{125}{\watt\per\meter\per\celsius}. The experimentally determined average aligns well within this range, confirming the accuracy of the method.

\subsection{Sources of Error}

\begin{itemize}
    \item \textbf{Thermocouple Calibration}: Any inaccuracies in thermocouple readings can affect temperature differences.
    \item \textbf{Contact Resistance}: Imperfect thermal contact between layers may introduce resistance to heat flow.
    \item \textbf{Assumption Limitations}: The assumption of one-dimensional heat flow may not hold perfectly due to edge effects.
\end{itemize}

\section{Conclusion}

The experiment successfully determined the thermal conductivity of brass, yielding an average value of \SI{121.06}{\watt\per\meter\per\celsius}. This result is consistent with established literature values, demonstrating the validity of the experimental approach and the simplifying assumptions made. While minor discrepancies and sources of error were identified, they did not significantly impact the overall findings. Future improvements could include better insulation and calibration of measurement devices to reduce experimental uncertainties.

\end{document}
